% 12_robot
\section{机器人拾球流程与本体侧优化}

网球拾球机器人是本报告内核侧工作的应用载体与验证对象。调度亲和、DVFS 与启动路径上的改动，最终都要落到这台机器人上电、追球、抓取、回桶的完整动作里接受检验。本体侧是一份用户态 C++ 可执行文件，主机 dry-run 与 RK3588 板端运行同一份状态机源码，仅替换电机与机械臂后端；因此感知链路的端到端时延在两处口径一致、可复现。感知在 480$\times$640 分辨率上的端到端时延\keyres{与 Linux 持平（\nInferTennis~ms 对 \nInferTennisLin~ms）}，上电到首帧推理（TTFI）领先 Linux（\nTtfi~s 对 \nTtfiLin~s），二者机理分别见推理链路一节与 TTFI 一节。把推理线程固定到 A76 大核簇、把解码交给 JPU、把色彩转换与缩放交给 RGA，是内核侧工作落到本体的直接抓手，其加速见推理链路与视觉感知两节。本节只讲本体侧的应用逻辑：纯逻辑拾球状态机、可换真实执行器后端，以及里程计引导的回桶。

运行时的数据流是一条闭环。应用从 UVC 相机取帧，感知按当前状态在两个检测器之间二选一：追球阶段把整理好的帧经 RGA 与 JPU 的零拷贝路径送入 NPU 上的单类 YOLOv8 检测头，DFL 解码后取置信度最高的球；寻桶阶段改用整型 HSV 红色掩膜，对满足色相与饱和度阈值的像素以四连通并查集求最大连通域来定位桶。该 YOLOv8 头会自动识别单类结构，其 DFL 解码此前只支持框通道数 64（\texttt{reg\_max} 为 16）的检测头，现已\ours{推广到框通道数 32（\texttt{reg\_max} 为 8）的紧凑头}，解码时按框通道数除以 4 推得 \texttt{reg\_max}，无需重训或更换后处理即可兼容更小的导出模型。感知结果进入状态机生成底盘与舵机指令，经去重后下发给差速底盘与 ZP10S 机械臂。\texttt{perception\_mode} 依状态发布感知模式，\texttt{CHASE\_BALL} 与 \texttt{GRAB} 走球检测器，其余状态走桶检测器，感知因此不会跑错检测器。以下按状态机、执行器后端、里程计回桶、工程加固四部分分述。

\subsection{拾球状态机}

\figph{../figures/out/F20.png}{拾球状态机：六状态非阻塞闭环，两个刹车时相由单调时钟推进，定时时相时长与相机供帧频率解耦}

拾球逻辑是一台六状态有限状态机，一轮完整循环沿 \texttt{CHASE\_BALL}、\texttt{GRAB}、\texttt{RETURN\_TO\_BUCKET} 或 \texttt{FIND\_BUCKET}、\texttt{APPROACH\_BUCKET}、\texttt{DEPOSIT} 推进，投放完成后复位里程计并回到 \texttt{CHASE\_BALL}。两个刹车时相 \texttt{BrakeForGrab} 与 \texttt{BrakeForDeposit} 穿插在抓取与投放之前，让底盘在动作前停稳。状态机的推进是纯逻辑的：单次 step 只依据当前检测与时钟给出下一步抽象指令，本身不做任何 I/O、不睡眠，串口写入与舵机阻塞等待都隔离在 dispatch 一侧。这样一来单帧感知无论快慢都不会阻塞状态推进，同一份状态机在主机 dry-run 与板端运行完全一致。

定时时相以单调时钟的绝对到期时刻计时。刹车保持 350~ms、抓取 settle 100~ms、释放 settle 500~ms 各记为一个 deadline，由每帧的 step 与帧间约 1~ms 一次的 tick 共同推进。因此定时时相的时长只取决于挂钟，与相机供帧频率解耦；相机偶发丢帧或帧率波动时，刹车与 settle 的时长保持不变。

\paragraph{追球对准}
\texttt{CHASE\_BALL} 采用固定速度的原地转向对准，按优先级从高到低逐条判定，先成立者即返回。其一，球丢失时朝上一次可见方向以固定速度原地转向搜索，转向由上一帧记下的横向偏移符号播种。其二，球过近（面积占比达到 \texttt{area\_reverse} 阈值 0.50）时直线倒车让开，倒车判定的优先级高于抓取，避免贴脸抓空。其三，横向偏移超过对准窗口 20 像素时以固定速度原地转向，把球对到目标列：夹爪在相机光心右侧约 90 像素，故目标列取列坐标 410，即在画面中心 320 之外右移 90 像素。其四，球尚远（面积占比低于 \texttt{area\_stop} 阈值 0.28）时直行趋近。其五，当面积占比落在 0.28 与 0.50 之间且偏移进入对准窗口，累计确认帧并输出刹车。追球全程只下发固定速度的离散可执行指令，此前的比例控制器与失速踢腿动作已删除，以适配只认固定速度的真实底盘。抓取触发经三帧确认防抖，连续三帧成立才启动 \texttt{BrakeForGrab} 刹车时相，刹车到期进入 \texttt{GRAB}。

\paragraph{抓取与两次脉冲确认}
\texttt{GRAB} 首次进入下发一次抓取动作并置 100~ms 的 settle deadline。抓取是否成功由\ours{两次脉冲阈值确认}：机械臂夹爪在贴地夹合与提起两个位姿各回读一次夹爪舵机的脉冲宽度，两次读数均达到 1260 微秒才判为夹住。其原理是球撑住夹爪、使其合不到空爪应有的脉冲值，故夹住时读数偏高。若判为空抓，状态机经 \texttt{on\_grab\_empty} 回到 \texttt{CHASE\_BALL} 重新搜索，避免空爪去桶。settle 到期后若里程计有效则进入 \texttt{RETURN\_TO\_BUCKET} 盲导回桶，否则进入 \texttt{FIND\_BUCKET} 靠视觉找桶。

\paragraph{找桶、趋近与投放}
\texttt{FIND\_BUCKET} 原地旋转搜索，以固定速度 12 转向；一旦检测到桶便累计确认帧，连续三帧成立后进入 \texttt{APPROACH\_BUCKET}。\texttt{APPROACH\_BUCKET} 是唯一使用比例转向的状态：偏置由桶心的列偏移按增益折算（等效于列偏移除以 16）并钳位到正负 8，死区 15 像素以内不转；前进速度在桶较远时取 25、面积占比达到 0.70 后降到 5 爬行。连续超过 10 帧丢失桶则退回 \texttt{FIND\_BUCKET}。桶的面积占比达到投放阈值 0.90 时启动 \texttt{BrakeForDeposit} 刹车时相，到期进入 \texttt{DEPOSIT}。\texttt{DEPOSIT} 首次进入下发一次释放动作并置 500~ms 的 settle deadline，到期后令机械臂复位、复位里程计锚点并回到 \texttt{CHASE\_BALL}，闭合一轮拾球。

\subsection{可换真实执行器后端}

\figph{../figures/out/F22.png}{执行器分层：\texttt{MotorBackend} 与 \texttt{ArmBackend} 接口下挂 Trace、PWM/DRV8833、ESP32-C3 UART 与 ZP10S UART 四个后端，主机 dry-run 锁定虚拟路径，真实后端初始化失败即致命}

状态机产生的是抽象指令，底盘一侧为 Drive、Brake、Standby，机械臂一侧为 Grab、Release、Ready；指令的物理实现被收敛到两个接口后面，即 \texttt{MotorBackend} 与 \texttt{ArmBackend}。一个字符串分派工厂按配置选定后端，共四种：底盘一侧有 Trace 虚拟、PWM/DRV8833 与 ESP32-C3 UART 三种，机械臂一侧有 Trace 虚拟与 ZP10S UART 两种。两种真实电机后端外层统一套一层死区装饰器，把任何非零指令抬升到电机的最小可动速度之上。本体应用此前只有 Trace 虚拟后端，把指令打成轨迹供主机基准与 CI 使用；现\ours{新增上述三个真实后端}，从纯虚拟基准变为可驱动真机，同时保留虚拟路径让端到端延迟保持纯计算、可复现。

同一份二进制既能在主机上以虚拟后端做 dry-run，又能在板上驱动真机。真实后端只在板端构建中编入，主机 dry-run 路径被锁定在虚拟后端；工厂在 ready 阶段一旦初始化失败即视为致命并终止本次运行，不回退虚拟后端，以免真机在无人察觉时空爪空转。部署由配置驱动，\texttt{-{}-config} 加载器先读配置文件再由命令行覆盖，一份二进制既能在 CI 做 dry-run，又能在板端部署，并对物理上不可执行的速度强校验拒绝。

\paragraph{ESP32-C3 底盘协议}
ESP32-C3 底盘走 115200 8N1 的 \texttt{/dev/ttyS6}。帧格式为帧头 0xAA、0x55 接命令字、载荷长度、载荷与异或校验，在线长度为 5 加载荷长度，接收侧先扫 0xAA 再扫 0x55 重同步。下发速度用命令 0x13，载荷为一对大端有符号百分比、钳位到正负 100，采用不等应答的即发即弃方式。里程计遥测另走命令 0x20，下位机以两帧 0x90 回报左右轮转速，供下一子节的里程计死推使用。

\paragraph{ZP10S 机械臂协议}
ZP10S 机械臂走 115200 的 \texttt{/dev/ttyS3}，使用阻塞式 ASCII 舵机协议。移动命令形如 \texttt{\#000P1611T1000!}，字段依次为舵机号、脉冲宽度与运动时间；位置回读命令形如 \texttt{\#002PRAD!}，回报的脉冲值正是上一子节抓取确认所依据的量。角度到脉冲的映射为 500 微秒加角度按 270 度满量程折算的分量，钳位到 500 至 2500 微秒。抓取被重写为一段标定位姿序列：张开、伸到位、合爪、回读贴地脉冲、抬起、回读提起脉冲，两次脉冲均达到 1260 微秒阈值才判为夹住。每步之间阻塞约 1~s，据此抓取、释放、复位分别约 5~s、3~s、1~s，这些阻塞落在 dispatch 一侧，不进入纯逻辑的状态推进。

\paragraph{PWM/DRV8833 底盘}
PWM/DRV8833 后端经 sysfs 驱动 H 桥，四个 pwmchip 分别对应左右轮的正反向输入，周期固定为 1~kHz，占空比按速度百分比线性折算。刹车令四路全部拉高，使 H 桥两端同高、主动短接电机；standby 令四路全关、任其滑行。右轮固定叠加正负 2 的偏置以补偿两侧电机差异。

\subsection{里程计引导回桶}

\figph{../figures/out/F21.png}{里程计引导回桶：投放时 \texttt{reset\_anchor} 把桶设为里程计原点，抓球后按到锚点的距离与方位盲导回桶，视觉一旦捕获桶立即抢占}

回桶依靠一套差速轮里程计。两轮的 RPM 遥测按差速模型死推位姿，轮半径取 0.03~m、轮距取 0.18~m，积分出机体的平面坐标与航向。投放动作完成时调用 \texttt{reset\_anchor} 把当前位姿清零，桶就此成为里程计原点；此后 \texttt{estimate} 随时给出到锚点的直线距离（坐标的欧氏模）与方位角。里程计默认关闭，仅在板端的 live 配置中开启。

抓球完成后若里程计有效，状态机进入 \texttt{RETURN\_TO\_BUCKET} 盲导回桶：按到锚点方位角与当前航向之差决定动作，航向误差超过 15 度时以固定速度原地转向收敛，误差在容差以内时以固定速度直行趋近，直到进入 0.50~m 的到达半径。回桶全程视觉具有最高优先级，一旦检测到桶立即抢占，转入 \texttt{FIND\_BUCKET} 交由视觉闭环趋近，桶位因此被记住而无需每轮重新搜索。超时 15~s、里程计失效或推算距离超出 10~m 的合理上限时，同样退回视觉搜索。若抓取完成时里程计无效，状态机直接进入 \texttt{FIND\_BUCKET} 靠视觉找桶。

为不让遥测阻塞控制，里程计采样被移出控制回路。一次同步的 RPM 读取需约 300~ms（一帧下发加两帧各 150~ms 超时回读），留在控制回路里会卡住电机决策。现\ours{改由一个后台线程按固定周期 100~ms 读取轮速}并更新位姿快照，控制回路每轮只取一次快照。同时对串口写入加锁、读取只清输入缓冲，以免遥测误丢尚未发出的驱动帧。慢速遥测与电机决策因此各行其道。

\subsection{工程加固与延迟工程}

\figph{../figures/out/F19.png}{机器人感知到控制的闭环：相机 latest-frame 锁存、感知二选一、状态机到 dispatch 的通路，无帧时经约 1~ms 的 tick 旁路推进定时时相}

现场运行要求这套闭环在异常下安全收场，本体侧为此叠加了若干加固。执行器接口的每个方法都返回成功与否，失败沿 dispatch 逐层冒泡，任一执行器调用失败即终止本次运行并令底盘 standby 停车。相机侧设两道保护：启动时的 warmup 要求连续三帧结构有效才进入主循环，超时上限 3000~ms；运行中的无帧看门狗在 2000~ms 内取不到新帧即停机退出。进程收尾由 RAII 接管 SIGINT 与 SIGTERM，每条退出路径都先令底盘 standby。下发侧对指令去重，相同的 Brake、Standby 或相同左右轮速的 Drive 不重复发送，减少串口流量与下位机抖动。推理侧以 \texttt{rknn\_set\_core\_mask} \ours{配置 NPU 的核掩码启用双核并发}，把检测吞吐摊到两个计算核上（第三核经消融确认无增益，详见推理链路一节）。

感知到控制的取帧路径也为低延迟设计。相机采集线程只把最新一帧锁存进单一互斥槽，消费者慢时永远取到最新帧、不积压旧帧。主循环有帧时跑感知加控制，无帧时每约 1~ms 打一次 tick、仅推进定时时相，控制节奏因此不被相机拖住。帧的时间戳在取帧时以单调时钟打点，未采用相机 PTS（该版本采集库不提供该字段），故端到端延迟统计少算了采集线程到取帧之间的排队，板端引用相关延迟时需注明这一口径。

本体侧的功能与加固逻辑均已由主机侧单元测试覆盖：状态机的时相边界精确到毫秒断言，执行器协议经伪终端断言，里程计积分经纯数学断言。就代码结构而言，机械臂的抓取、释放、复位分别由五步、三步、一步逐步阻塞的位姿序列组成，每步之间阻塞约 1~s，据此给出抓取约 5~s、释放约 3~s、复位约 1~s 的阻塞下界。
